\section{Proof of the Main Theorem}

\begin{proof}[Proof of Theorem 2.1]
Proposition \ref{prop:scale-barrier} suppresses the dangerous high-frequency cascade on the classical equation itself. Proposition \ref{prop:compactness} closes the quadratic nonlinear term on the same classical approximation surface, so no defect survives in the convective term. Proposition \ref{prop:gradient-control} then converts the same classical dynamics into a Euclidean continuation inequality that prevents derivative-scale blow-up.

These three statements live on one admissibility surface: same equation, same data class, same viscosity regime, and same pressure law. No theorem-critical step is borrowed from a modified equation, geometric damping mechanism, or stronger hidden hypothesis. The only singularity route tracked in the present framework is therefore blocked on the classical surface itself. Hence the conclusion of Theorem \ref{thm:classical-closure} holds.
\end{proof}

\begin{remark}[Proof architecture]
The proof closes by the interaction of the three bridge propositions, not by any one of them in isolation. The scale barrier supplies the dangerous-scale suppression used by compactness and gradient control, the compactness package closes the nonlinear term on the same equation, and the gradient-transfer inequality prevents derivative-scale re-entry of singular behavior.
\end{remark}
